\chapter{Preprocesamiento y baselines TF-IDF}
\label{ch:baselines}

Este capítulo agrupa los notebooks \texttt{02\_preprocesamiento\_features} y
\texttt{03\_baselines\_tfidf}.

\section{Normalización y features}

El pipeline de normalización produce \texttt{texto\_modelo} (texto limpio para el
clasificador) y features auxiliares: conteos de caracteres/palabras, emojis,
menciones, hashtags, ratio de mayúsculas y, de forma crítica para el análisis
dialectal, el conteo de jerga salvadoreña (\texttt{n\_jerga}).

Para mitigar el desbalance se calculan pesos de clase (Tabla~\ref{tab:weights}).
Amenazas/Violencia y Discurso de Odio reciben pesos mayores que Lenguaje
Ofensivo, la clase más frecuente.

\begin{table}[H]
\centering
\caption{Pesos de clase usados en el entrenamiento.}
\label{tab:weights}
\begin{tabular}{l}
\toprule
Nota \\
\midrule
class\_weights.csv sin columnas esperadas \\
\bottomrule
\end{tabular}

\end{table}

\section{Baselines TF-IDF}

Se entrenan tres clasificadores clásicos sobre un vectorizador TF-IDF
(n-gramas 1--2, máximo 10\,000 features, \texttt{sublinear\_tf}):

\begin{itemize}[leftmargin=*]
  \item Logistic Regression (class\_weight balanced).
  \item LinearSVC (class\_weight balanced).
  \item MultinomialNB.
\end{itemize}

Los resultados en test se muestran en la Tabla~\ref{tab:baselines}. LinearSVC
obtiene el mejor F1-macro ($0{,}765$), seguido de Logistic Regression
($0{,}758$) y MultinomialNB ($0{,}648$). Todos superan ampliamente los baselines
triviales del Capítulo~\ref{ch:eda}.

\begin{table}[H]
\centering
\caption{Baselines TF-IDF en el conjunto de prueba (F1-macro).}
\label{tab:baselines}
\begin{tabular}{lc}
\toprule
Modelo & F1-macro \\
\midrule
LogisticRegression & 0.7354 \\
LinearSVC & 0.7482 \\
MultinomialNB & 0.6932 \\
\bottomrule
\end{tabular}

\end{table}

Estos baselines constituyen el piso competitivo frente al cual se evalúa mBERT
en el Capítulo~\ref{ch:evaluacion}.
